\documentclass[paper]{ieice}
\usepackage{float}
\usepackage{cite}
\usepackage{caption}
\usepackage{subcaption}
\usepackage{amsmath,amssymb,amsfonts}
% \usepackage{algorithmic}
\usepackage[dvipdfmx]{graphicx}
\usepackage{textcomp}
\usepackage{xcolor}
\usepackage{colortbl}
\usepackage[T1]{fontenc}
\usepackage{mdframed}
\usepackage{algorithm, algpseudocode} % texlive-science
\usepackage[braket, qm]{qcircuit}


\usepackage[braket, qm]{qcircuit}
\usepackage{tikz}
%
\newcommand*\rtofstarg{\tikz[baseline=(char.base)]{
            \node[shape=circle,double,draw,inner sep=1pt,minimum width=16pt] (char) {S};}}
\newcommand*\rtofsdgtarg{\tikz[baseline=(char.base)]{
            \node[shape=circle,double,draw,inner sep=1pt] (char) {S$^{\dagger}$};}}
\newcommand*\rtoftarg{\tikz[baseline=(char.base)]{
            \node[shape=circle,draw,inner sep=-0.5pt,fill=none] (char) {$\oplus$};}}
\newcommand*\rtofsctrl{\tikz[baseline=(char.base)]{
            \node[shape=circle,double,draw,inner sep=1pt,fill=black, minimum width=3pt] (char) {};}}
\newcommand*\rtofsctrlo{\tikz[baseline=(char.base)]{
            \node[shape=circle,double,draw,inner sep=1pt,fill=white, minimum width=3pt] (char) {};}}

\newcommand*\ctrlrtofs[1]{\push{\rtofsctrl}\qwx[#1]\qw}
\newcommand*\ctrlortofs[1]{\push{\rtofsctrlo}\qwx[#1]\qw}
\newcommand*\targrtofs{\push{\rtofstarg}\qw}
\newcommand*\targrtof{\push{\rtoftarg}\qw}


\algrenewcommand\algorithmicindent{0.5em}%
\newtheorem{definition}{Definition}
\newtheorem{example}{Example}
\newcommand{\thickhline}{%
    \noalign {\ifnum 0=`}\fi \hrule height 1pt
    \futurelet \reserved@a \@xhline
}
\def\BibTeX{{\rm B\kern-.05em{\sc i\kern-.025em b}\kern-.08em
    T\kern-.1667em\lower.7ex\hbox{E}\kern-.125emX}}
\newcommand{\figcaption}[1]{\def\@captype{figure}\caption{#1}}
\newcommand{\tblcaption}[1]{\def\@captype{table}\caption{#1}}

\newcount\NGCcolor
\NGCcolor=1
%\NGCcolor=0
\ifnum\NGCcolor<1
  \newcommand{\redout}[1]{\textcolor{red}{\bf #1}}
  \newcommand{\blueout}[1]{\textcolor{blue}{\bf #1}}
  \newcommand{\greenout}[1]{\textcolor{green}{\bf #1}}
  \newcommand{\cyannout}[1]{\textcolor{cyan}{\bf #1}}
  \newcommand{\magentaout}[1]{\textcolor{magenta}{\bf #1}}
  \newcommand{\yellowout}[1]{\textcolor{yellow}{\bf #1}}
\else
  \newcommand{\redout}[1]{#1}
  \newcommand{\blueout}[1]{#1}
  \newcommand{\greenout}[1]{#1}
  \newcommand{\cyannout}[1]{#1}
  \newcommand{\magentaout}[1]{#1}
  \newcommand{\yellowout}[1]{#1}
\fi

\setcounter{page}{1}

\field{A}
\title{Leveraging Different Boolean Function Decompositions to Optimize LUT-based Quantum Circuit Synthesis}
\authorlist{%
 \authorentry[dizzy@ngc.is.ritsumei.ac.jp]{David Clarino}{n}{rits}\MembershipNumber{}
 \authorentry[accel@ngc.is.ritsumei.ac.jp]{Naoya Asada}{n}{rits}\MembershipNumber{}
 \authorentry[ger@cs.ritsumei.ac.jp]{Shigeru Yamashita}{e}{rits}\MembershipNumber{}
}
\affiliate[rits]{The author is with the Graduate School of Information Science and Engineering, Ritsumeikan University, Shiga, 525-8577 Japan.
}


\begin{document}
\maketitle

\begin{summary}
  LUT-based synthesis methods have recently been proposed as a way to
  synthesize quantum Boolean circuits in a qubit constrained
  environment. Other recent research has demonstrated the possibility of using
  Relative Phase quantum circuits in tandem has an advantage in minimizing
  T-count in the fault-tolerant quantum computing paradigm. Previously a Shannon's
  Decomposition based algorithm for generating Relative Phase quantum circuits,
  used in conjunction with LUT-based synthesis, yields a small but clear advantage
  in T-counts. In order to further this advantage, several other techniques
  are integrated, such as using Davio expansions, as well as alternating
  balanced and unbalanced Relative Phase Circuits. This yields a much more clear
  advantage even at bigger LUT-sizes. 
  
\end{summary}
\begin{keywords}
relative-phase Toffoli gates (RTOF), circuit optimization, LUT, synthesis
\end{keywords}
\setlength\textfloatsep{4pt}
\setlength\floatsep{0pt}
\setlength\abovecaptionskip{4pt}


\begin{figure*}[t]
    \centering
    \begin{minipage}{0.6\textwidth}
        \begin{subfigure}[b][15mm]{0.9\textwidth}
            \centering
             \Qcircuit @C=0.5em @R=0.2em @!R { \\
          \lstick{{x}_{0}} & \ctrl{2}   & \qw &          & & \push{\ket{x_1}}  & & \qw               & \qw      & \qw                       & \ctrl{2} & \qw               & \qw      & \qw                       & \ctrl{2} & \ctrl{1}          & \gate{\mathrm{T}}         & \ctrl{1} & \qw & \qw & & \ket{x_1}\\
          \lstick{{x}_{1}} & \ctrl{1}   & \qw & \push{=} & & \push{\ket{x_2}}  & & \qw               & \ctrl{1} & \qw                       & \qw      & \qw               & \ctrl{1} & \gate{\mathrm{T}}         & \qw      & \targ             & \gate{\mathrm{T^\dagger}} & \targ    & \qw & \qw & & \ket{x_2}\\
          \lstick{{f}}     & \targ \qw  & \qw &          & & \push{\ket{x_3}} & & \gate{\mathrm{H}} & \targ    & \gate{\mathrm{T^\dagger}} & \targ    & \gate{\mathrm{T}} & \targ    & \gate{\mathrm{T^\dagger}} & \targ    & \gate{\mathrm{T}} & \gate{\mathrm{H}}         & \qw      & \qw & \qw & & \push{\ket{x_1}\ket{x_1 \cdot x_2 \oplus x_3}}\\
}

            \caption{A Toffoli gate in Clifford+T basis}
            \label{fig-toffoli}
        \end{subfigure}
        
        \vspace{10mm}
        \begin{subfigure}[b][15mm]{0.9\textwidth}
            \centering
             \Qcircuit @C=0.5em @R=0.2em @!R { \\
          \lstick{{x}_{0}} & \ctrl{2}\qw         & \qw &          & & \push{\ket{x_1}} & & \push{{x}_{0} :  }  & \qw               & \qw                       & \qw      & \qw                       & \ctrl{2} & \qw               & \qw      & \qw                         & \qw               & \qw & \\
          \lstick{{x}_{1}} & \ctrl{1}            & \qw & \push{=} & & \push{\ket{x_2}} & & \push{{x}_{1} :  }  & \qw               & \qw                       & \ctrl{1} & \qw                       & \qw      & \qw               & \ctrl{1} & \qw                         & \qw               & \qw & \\
          \lstick{{f}}     & \push{\rtoftarg}\qw & \qw &          & & \push{\ket{x_3}} & &  \push{{f}     :  } & \gate{\mathrm{H}} & \gate{\mathrm{T^\dagger}} & \targ    & \gate{\mathrm{T^\dagger}} & \targ    & \gate{\mathrm{T}} & \targ    & \gate{\mathrm{T^{\dagger}}} & \gate{\mathrm{H}} & \qw & \\
}

            \caption{An RTOF gate in Clifford+T basis}
            \label{fig-rtof}
        \end{subfigure}
    \end{minipage}
    \hspace{10mm}
    \begin{minipage}{0.3\textwidth}
    \vspace{10mm}
        \begin{subfigure}[b][30mm]{\textwidth}
        \centering
        \begin{tabular}{c|c}      
            \hline
            {x1,x2,x3} & $\Delta\theta$  \\
            \hline
            000   & 0      \\
            001   & 0      \\
            010   & 0      \\
            011   & 0      \\
            100   & 0      \\
            101   & $+\pi$          \\
            110   & $+\pi/2$   \\
            111   & $-\pi/2$ \\
            \hline
        \end{tabular}
        \caption{Phases added by an RTOF gate}
        \label{table-rtof}
        \end{subfigure}
    \end{minipage}
    \caption{A Toffoli gate and an RTOF gate}
    \label{fig-tof-rtof}
\end{figure*}

\section{Introduction}

{\it Quantum Boolean circuits}~\cite{bib-ddmf} are common components of many quantum algorithms such as oracles in Grover's Algorithm and cost functions in VQE. To realize such quantum Boolean circuits, a Toffoli gate is an essential logic primitive in a universal gate set. These Toffoli gates have to be, in turn, composed of physically realizable gates, such as NOT, CNOT, and T. (e.g., the Clifford+T~\cite{bib-barenco-elementary} basis gate set), as depicted in Fig.~\ref{fig-toffoli}. However, in the fault-tolerant paradigm, T-gates incur much higher cost than so-called Clifford gate types. It is therefore useful to consider ways to reduce this {\it T-count}.

One of the methods that has been proposed to reduce this T-count is to implement Boolean functions only up to a {\it relative phase}~\cite{bib-barenco-elementary}. Some recent work~\cite{bib-rtof-maslov}~\cite{bib-selinger-rtof}~\cite{bib-amy-phase-state} has proposed relative phase versions of common Boolean functions, particularly the Relative Phase Toffoli Gate (RTOF). 
The work in~\cite{bib-rtof-maslov} showed that the method of constructing circuits offers advantages in T-count.

Both~\cite{bib-rtof-maslov}~and~\cite{bib-amy-phase-state} have also proposed methods to generalize the construction of relative phase versions of quantum circuits that were introduced in~\cite{bib-selinger-rtof}. Because of these findings, cascading logic, such as those produced by LUT-based synthesis methods, thus present an interesting candidate to integrate with relative phase methods.

Recent attempts to use~\cite{bib-lhrs-soeken} to synthesize quantum circuits using LUT-based methods have shown some promising results. As intermediate nodes ``cascade'' their outputs into output nodes, logic produced by such a method is exactly the type that can be improved by synthesizing relative phase versions of its component logic. In addition, the fact that relative phase versions can be implemented without any additional ancilla means that it works well to synthesize Boolean functions in the qubit constrained environment that such methods are intended for. 

Our research seeks to leverage recent results in relative phase implementations of Boolean functions and LUT-based quantum circuit synthesis to reduce T-count in a quantum circuit.

\noindent {\bf Our Contribution.} 
We propose a method that reduces the number of T-gates in a LUT-based synthesized quantum Boolean circuit by improving upon our existing method \footnote{The preliminary version of this research will appear in ICCAD} to generate relative phase quantum circuits. Among our contributions:
\begin{itemize}
\item A method to generate arbitrary Boolean functions up to a relative phase using Shannon's Decomposition and Davio Expansions, as well as balanced and unbalanced constructions of relative phase gates.
\item A method to generate a quantum Boolean circuit from a Boolean logic network representation leveraging the above method to reduce T-count.
\end{itemize}

The paper is structured as follows: Section~\ref{Sec:prelim} outlines preliminary information for understanding this paper. Section~\ref{Sec:proposed} follows with our proposed method, including both the method to realize functions up to a relative phase without any additional ancilla, as well as its use with the LUT-based synthesis. Then, our experimental results in Section~\ref{Sec:exp} compares the proposed method to the existing method and the naive method and finds a clear advantage in its use. Finally Section~\ref{Sec:concl} contains our conclusions and future work.




% \begin{figure*}[ht]
  
%   \begin{minipage}{.46\linewidth}
%     \begin{center}
%       \Qcircuit @C=0.5em @R=0.2em @!R { \\
          \push{{x}_{0} }   & & \multigate{3}{x_k \cdot U}  & \qw   & \qw & &         & & \qw & \qw               & \qw                       & \qw      & \qw                       & \multigate{3}{U}              & \qw               & \qw      & \qw                         & \qw               & \qw & \\
          \push{{x}_{1} }   & & \ghost{x_k \cdot f}         & \qw   & \qw & &         & & \qw & \qw               & \qw                       & \qw      & \qw                       & \ghost{U}                      & \qw               & \qw      & \qw                         & \qw               & \qw & \\
          \push{\vdots}     & &                             &       &     & &\push{=} & &     &                   & \push{\vdots}             &          &  \push{\vdots}            &                                & \push{\vdots}     &          &                             & \push{\vdots}     &     & \\
          \push{{x}_{n-1} } & & \ghost{x_k \cdot f}         & \qw   & \qw & &         & & \qw & \qw               & \qw                       & \ctrl{1} & \qw                       & \ghost{U}                      & \qw               & \ctrl{1} & \qw                         & \qw               & \qw & \\
          \push{{x}_{n} }   & & \push{\rtoftarg}\qw\qwx[-1] & \qw   & \qw & &         & & \qw & \gate{\mathrm{H}} & \gate{\mathrm{T}}         & \targ    & \gate{\mathrm{T^\dagger}} & \targ\qwx[-1]    & \gate{\mathrm{T}} & \targ    & \gate{\mathrm{T^{\dagger}}} & \gate{\mathrm{H}} & \qw & \qw\\
}
                            

%     \end{center}
%     \caption{U-controlled RTOF gate}
%     \label{fig-urtof}
%   \end{minipage}
%   %\vspace*{0.25cm}
% \end{figure*}
% \vspace*{0.1cm}
\section{Preliminary}
\label{Sec:prelim}
\subsection{Quantum Boolean Circuits}

In this paper, we consider {\it quantum Boolean circuits}~\cite{bib-ddmf} that implement Boolean functions. We will be implementing them using the fault-tolerant set of Clifford gates, T gates, and S gates. 

We denote the primary inputs of a circuit by $x_1, x_2, \cdots,x_{n}$  when we have $n$ primary inputs. Each qubit state in a circuit can be represented by a Boolean function with respect to $x_1, x_2, \cdots,x_{n}$.

An example of such a quantum Boolean operation is the Toffoli gate depicted in
Fig ~\ref{fig-toffoli}. It inverts the basis state of the target bit (i.e.,
$\ket{1}$ to $\ket{0}$ and $\ket{0}$ to $\ket{1}$) when the states of the both
control bits are $\ket{11}$. This provides us a a way to realize the AND
operation in a quantum circuit.

\subsection{Relative-Phase Boolean Functions}

A typical decomposition of the Toffoli gate into the Clifford + T basis
consists of seven T gates as shown in Fig.~\ref{fig-toffoli}.
A Relative-Phase Toffoli (RTOF) gate~\cite{bib-rtof-maslov}, in contrast,
implements the same function using only four T gates as shown in
Fig.~\ref{fig-rtof}. However, it does this at the cost of a {\it relative phase},
which is an input dependent phase on the output.
Note the double circle on the target of Fig.~\ref{fig-rtof}: we will be
using this notation to denote a relative phase version of a logic
gate in this paper. One may notice that the version of the RTOF gate
that is used here is similar to the $n$-controlled-iX gate
from~\cite{bib-selinger-rtof} for $n=2$, with one less CNOT gate. As
noted in~\cite{bib-rtof-maslov}, removing this gate only affects the
phase, so it only transforms an RTOF gate into another RTOF gate with
the same T-count. We will be considering analogs of both types of
relative phase gatesin Section ~\ref{Sec:proposed}

As depicted in Fig.~\ref{fig-rtof}, an RTOF gate has three inputs; the input quantum basis states are $\ket{000}$
to $\ket{111}$. Similar to Toffoli gates, an RTOF gate swaps  $\ket{110}$ and $\ket{111}$. In addition to this
logic operation, RTOF gates add phases to some output states that the equivalent Toffoli gate does not. The added
phases are as shown in Fig.~\ref{table-rtof}. For example, when the input state is $\ket{110}$, a phase $+\pi/2$ is
added.

One of the results from~\cite{bib-rtof-maslov} is that when larger functions are broken down into compute/uncompute pairs, such as in the Barenco style Toffoli decompositions, it is almost always more advantageous to implement each ``half'' of the compute/uncompute pair into a relative phase construction. We will next go into one such construction.
%\vspace{-0.25cm}
\subsection{Boolean Logic Network}

A {\it Boolean logic network} is a simple Directed Acyclic Graph (DAG) whose vertices are primary inputs, primary outputs, and nodes which represent logic gates, whose arcs connect gates to inputs, outputs, and other gates.

A node has zero or more {\it fanins}, or incoming edges, and zero or more {\it fanouts} i.e. outgoing edges. The {\it primary inputs (PIs)} are nodes without fanins in the current network. The {\it primary outputs} are a subset of nodes of the network without fanouts. 

A logic function can be decomposed into such a Boolean logic network, and its nodes subsequently implemented as LUTs or logic gates. We see an example in Fig.~\ref{fig-lut-network}.

\begin{figure}[tbp]
    \centering
    \usetikzlibrary{positioning}
\definecolor{offwhite}{HTML}{F2EDED}
%% \tikzset{
%%     > = stealth,
%%     every node/.append style = {
%%         draw = none,
%%         text = offwhite
%%     },
%%     every path/.append style = {
%%         arrows = ->,
%%         draw = offwhite,
%%         fill = offwhite
%%     },
%%     hidden/.style = {
%%         draw = offwhite,
%%         shape = circle,
%%         inner sep = 1pt
%%     }
%% }
%% \tikz{
%%     \node (x) {$X$};
%%     \node[hidden] (z) [above right = of x] {$Z$};
%%     \node (y) [below right = of z] {$Y$};
%%     \path (x) edge (y);
%%     \path (z) edge (x);
%%     \path (z) edge (y);
%% }
\tikzset{
  > = stealth,
  every node/.append style = {
    draw      = black,
    shape     = circle,
    inner sep = 1
  },
  every path/.append style = {
        arrows = ->,
        draw = black,
        fill = offwhite
  },  
  input/.style = {
    draw = none
  }
}

\tikz{
  \node[input] (x1) at (0,0) {$x1$};
  \node[input] (x2) at (1,0) {$x2$};
  \node[input] (x3) at (2,0) {$x3$};
  \node[input] (x4) at (3,0) {$x4$};
  \node[input] (x5) at (4,0) {$x5$};
  \node[input] (y1) at (1,3) {$y1$};
  \node[input] (y2) at (3,4) {$y2$};
  \node (1) at (1,1) {$1$};
  \node (2) at (3,1) {$2$};
  \node (3) at (1,2) {$3$};
  \node (4) at (3,2) {$4$};
  \node (5) at (3,3) {$5$};      
  \path (x1) edge (3);
  \path (x2) edge (1);
  \path (x3) edge (1);
  \path (x4) edge (1);
  \path (x4) edge (2);
  \path (x5) edge (2);
  \path (x5) edge (5);
  \path (1)  edge (3);
  \path (1)  edge (4);    
  \path (2)  edge (4);
  \path (4)  edge (5);
  \path (3)  edge (y1);
  \path (5)  edge (y2);  
}

    \caption{A Boolean logic network}
    \label{fig-lut-network}
\end{figure}


\section{Existing Research and Motivation}
 \label{Sec:Exist}

\subsection{Using RTOF Gates to Decompose MCT Gates}
Existing research~\cite{bib-rtof-maslov}~\cite{bib-selinger-rtof} finds an advantage to using RTOF gates in synthesizing and decomposing logic functions, particularly Multiply-Controlled Toffoli Gates (MCTs).

In particular, Maslov proposes the decomposition depicted in Fig.~\ref{fig-maslov-decomp}~\cite{bib-rtof-maslov}. Here $U$ is a unitary operation implemented up to a relative phase.

Maslov further posits a way to construct $n$-input RTOF gates by replacing the middle CNOT gate in Fig.~\ref{fig-rtof} with an $n$-input RTOF gate, and then repeating recursively. However, as Amy et. al.~\cite{bib-amy-phase-state} note, the form of this RTOF gate generates a target-dependent phase which, when combined with the Hadamard gate at the end, adds the terms of its phase function to the Boolean function to the output, such that it takes the form $F = x_k \cdot F_{k-1} + F_{k-2}$.


\begin{figure}[tbp]
    \centering
    \scalebox{1}{\Qcircuit @C=0.5em @R=0.2em @!R { \\
  \push{\ket{{x}_{0}} } & & \multigate{4}{F} & \qw & &   & & \multigate{2}{F'} & \qw & \qw           & \qw & \multigate{2}{{F'}^{\dagger}} & \qw & \rstick{\ket{{x}_{0}}} \\
  \push{\ket{{x}_{1}} } & & \ghost{F}        & \qw & &   & & \ghost{F'}        & \qw & \qw           & \qw & \ghost{{F'}^{\dagger}}        & \qw & \rstick{\ket{{x}_{1}}}  \\
  \push{\ket{{x}_{2}} } & & \ghost{F}        & \qw & &   & & \ghost{F'}\qwx[1] & \qw &               & \qw & \ghost{{F'}^{\dagger}}\qwx[1] & \qw & \rstick{\ket{{x}_{2}}}  \\
  \push{\vdots}         & &                  &     & & = & & \push{\vdots}     &     & \push{\vdots} &     & \push{\vdots}                 &     & \rstick{\vdots }    \\       
  \push{\ket{{x}_{k}} } & & \ghost{F}        & \qw & &   & &  \qw              & \qw & \ctrl{2}      & \qw & \qw                           & \qw & \rstick{\ket{{x}_{k} }} \\
  \push{\ket{0} }       & & \qw              & \qw & &   & & \targ\qwx[-2]     & \qw & \control\qw   & \qw & \targ\qwx[-2]                 & \qw & \rstick{\ket{0} }  \\    
  \push{\ket{0} }       & & \targ\qwx[-2]    & \qw & &   & & \qw               & \qw & \targ         & \qw & \qw                           & \qw & \rstick{\ket{F} } \\
}

}
    \caption{Toffoli gate decomposition from~\cite{bib-rtof-maslov}}
    \label{fig-maslov-decomp}
\end{figure}


\begin{figure}[tbp]
    \centering
    \scalebox{1}{\Qcircuit @C=0.5em @R=0.2em @!R { \\
  \lstick{\ket{{x}_{1}}} & \qw  & \qw                 & \qw                 & \multigate{5}{U_3}      & \qw                & \qw                   & \qw                          & \qw                          & \qw                           & \ghost{x_1}  & \lstick{\ket{{x}_{1}}} \\
  \lstick{\ket{{x}_{2}}} & \qw  & \multigate{1}{U_1}  & \qw                 & \push{\dashedwire}\qw   & \qw                & \qw                   & \qw                          & \qw                          & \multigate{1}{U_1^{\dagger}}  & \ghost{x_2}  & \lstick{\ket{{x}_{2}}} \\
  \lstick{\ket{{x}_{3}}} & \qw  & \ghost{1}           & \qw                 & \push{\dashedwire}\qw   & \qw                & \qw                   & \qw                          & \qw                          & \ghost{U_1^{\dagger}}         & \ghost{x_3}  & \lstick{\ket{{x}_{3}}} \\
  \lstick{\ket{{x}_{4}}} & \qw  & \qw                 & \multigate{1}{U_2}  & \push{\dashedwire}\qw   & \qw                & \qw                   & \qw                          & \multigate{1}{U_2^{\dagger}} & \qw                           & \ghost{x_4}  & \lstick{\ket{{x}_{4}}} \\
  \lstick{\ket{{x}_{5}}} & \qw  & \qw                 & \ghost{2}           & \push{\dashedwire}\qw   & \qw                & \multigate{4}{U_5}    & \qw                          & \ghost{U_2^{\dagger}}        & \qw                           & \ghost{x_5}  & \lstick{\ket{{x}_{5}}} \\
  \lstick{\ket{0}     }  & \qw  & \targ\qwx[-3]\qw    & \qw                 & \ghost{3}\qw            & \multigate{1}{U_4} & \push{\dashedwire}\qw & \multigate{1}{U_4^{\dagger}} & \qw                          & \targ\qwx[-3]\qw              & \ghost{0}    & \lstick{\ket{0}     }    \\
  \lstick{\ket{0}     }  & \qw  & \qw                 & \targ\qwx[-2]\qw    & \qw                     & \ghost{4}          & \push{\dashedwire}\qw & \ghost{U_4^{\dagger}}        & \targ\qwx[-2]\qw             & \qw                           &  \qw            & \lstick{\ket{0}     }  \\
  \lstick{\ket{0}     }  & \qw  & \qw                 & \qw                 & \targ\qwx[-2]\qw        & \qw                & \push{\dashedwire}\qw & \qw                          & \qw                          & \qw                           & \ghost{y_1}  & \lstick{\ket{y_1}   }  \\      
  \lstick{\ket{0}     }  & \qw  & \qw                 & \qw                 & \qw                     & \targ\qwx[-2]\qw   & \ghost{5}             & \targ\qwx[-2]\qw             & \qw                          & \qw                           & \ghost{0}    & \lstick{\ket{0}     }    \\
  \lstick{\ket{0}     }  & \qw  & \qw                 & \qw                 & \qw                     & \qw                & \targ\qwx[-1]\qw      & \qw                          & \qw                          & \qw                           & \ghost{y_2}  & \lstick{\ket{y_2}     }  \\
}
}
    \caption{The quantum circuit implementing the Boolean logic network in Fig.~\ref{fig-lut-network} using the LHRS. }
    \label{fig-lut-qc}
\end{figure}

To summarize:
\begin{itemize}
\item Replacing pairs of computing and uncomputing gates with RTOF gate and its inverse will lead to T gate savings in most cases.
\item Using the above, it is possible to use clean ancilla to decompose an MCT with fewer T gates than the known methods.
\item By utilizing a relative UCNOT implementation, we get a recursive procedure to generate MCT gates up to a relative phase.
\end{itemize}



\subsection{LUT-based Hierarchical Reversible Logic Synthesis (LHRS)}
\label{Sec:lhrs}

Previous research~\cite{bib-lhrs-soeken} has detailed a way to leverage LUT-based synthesis of Boolean functions to create reversible logic networks that will be essential to quantum Boolean circuits. The overall flow maps each node in the Boolean network to a function-controlled NOT gate (UCNOT), a NOT gate controlled not just by the value of a variable but that of a Boolean function, and then cascades them together to create the output. 

The result of the method is shown in Fig~\ref{fig-lut-qc}. To get from Fig.~\ref{fig-lut-network} to Fig~\ref{fig-lut-qc} first, we take as input a LUT network synthesized from a Boolean expression using existing LUT-Network decomposition methods. Then this LUT network is parsed by the algorithm to create a quantum circuit of UCNOT gates. Any intermediate node is decomposed as two UCNOTs: one to compute, and the other to uncompute, acting on the same dirty ancilla. Any node that acts only on output qubits (output node i.e., an endpoint to the cascading logic) is given only one UCNOT.


Because the UCNOT gates act on ancilla, there is a need to uncompute them, so the flow next produces uncompute logic. This produces the compute/uncompute pair that a relative-phase based construction of the function can take advantage of. 

\begin{example}
  \label{ex-lut}
  We start at node 1. Since this is an intermediate node, we place a
  UCNOT on an ancilla, and place its inverse on the same ancilla to
  uncompute. We do the same for 2. For node 4, we cascade the ancilla
  from nodes 1 and 2 into it as input, and place its LUT, as well as
  its uncompute logic, on another ancilla. Finally, nodes 3 and 5 are
  output nodes, so we only generate their compute logic. 

\end{example}


%%% Local Variables:
%%% mode: japanese-latex
%%% TeX-master: "main"
%%% End:

%% \begin{figure}
%%   \label{fig-maslov-decomp}
%% \end{figure}


%% \section{Existing Research and Motivation}

%% Because not all Boolean functions can be implemented as unitaries, from the very early days of Quantum Computing~\cite{bib-barenco-elementary}, it has been known that in addition to adding ancilla to the circuit to store temporary values, allowing the final states to have differing phases ({\it relative phase}) with respect to each other is also a valid way to implementa Boolean function.

%% Relative phase versions of quantum operations, such as the {\it Relative-phase Toffoli gates (RTOF)}~\cite{bib-rtof-maslov} ~\cite{bib-selinger-rtof} have been proposed as an alternative . One of the proposed implementations of the RTOF depicted in Fig.~\ref{fig-rtof} can calculate the same logic function as a Toffoli gate and reduces the T-gate count by 3. However it introduces a {\it relative phase} among its output quantum states, (enumerated in Table ~\ref{table-rtof}. Because of the introduced relative phase, we cannot simply replace a Toffoli gate with an RTOF in general. Rather, RTOFs can be utilized when we decompose a multiple-control Toffoli gate because the relative phase errors can be canceled between two RTOFs in the decomposition ~\cite{bib-rtof-maslov}.

%% Maslov further posits that not just Toffoli gates but generalized controlled U functions can take advantage of the fact that Boolean functions implemented up to a relative phase used in tandem can optimize the overall circuit.

%% Meanwhile, LUT-based quantum circuit synthesis has been proposed as a way to synthesize quantum circuits in an environment where available clean ancillae for use as intermediate qubits is limited. Methods such as the LHRS ~\cite{bib-lhrs-soeken}  synthesize each of these LUTs is implemented as a U-controlled NOT function. Naturally, this leads to a cascading series of gates, along with their uncompute pairs. This structure naturally lends itself to the construction from Maslov, where alternating pairs of gates can cancel their relative phases.

%% In light of these facts, it is easy to imagine a method that utilizes relative phase Boolean functions to reduce the cost of gates in the LUT network, all while adding no additional ancillae. The key will be a method to generate relative phase Quantum Boolean Circuits from arbitrary Boolean functions. We propose such a method in this work.





%% \begin{figure*}[t]
%%   \label{fig-cx-star}
%%   \begin{center}
%%   \end{center}
%%   \caption{The $C-X^{*}$ gate from \cite{bib-amy-phase-state}}
%% \end{figure*}
%% \begin{figure*}[t]
%%   \begin{minipage}{.46\linewidth}[t]
%%     \label{fig-lut-network}
%%     \begin{center}
%%       \Qcircuit @C=0.5em @R=0.2em @!R { \\
  \lstick{\ket{{x}_{1}}} & \qw  & \qw                 & \qw                 & \multigate{5}{U_3}      & \qw                & \qw                   & \qw                          & \qw                          & \qw                           & \ghost{x_1}  & \lstick{\ket{{x}_{1}}} \\
  \lstick{\ket{{x}_{2}}} & \qw  & \multigate{1}{U_1}  & \qw                 & \push{\dashedwire}\qw   & \qw                & \qw                   & \qw                          & \qw                          & \multigate{1}{U_1^{\dagger}}  & \ghost{x_2}  & \lstick{\ket{{x}_{2}}} \\
  \lstick{\ket{{x}_{3}}} & \qw  & \ghost{1}           & \qw                 & \push{\dashedwire}\qw   & \qw                & \qw                   & \qw                          & \qw                          & \ghost{U_1^{\dagger}}         & \ghost{x_3}  & \lstick{\ket{{x}_{3}}} \\
  \lstick{\ket{{x}_{4}}} & \qw  & \qw                 & \multigate{1}{U_2}  & \push{\dashedwire}\qw   & \qw                & \qw                   & \qw                          & \multigate{1}{U_2^{\dagger}} & \qw                           & \ghost{x_4}  & \lstick{\ket{{x}_{4}}} \\
  \lstick{\ket{{x}_{5}}} & \qw  & \qw                 & \ghost{2}           & \push{\dashedwire}\qw   & \qw                & \multigate{4}{U_5}    & \qw                          & \ghost{U_2^{\dagger}}        & \qw                           & \ghost{x_5}  & \lstick{\ket{{x}_{5}}} \\
  \lstick{\ket{0}     }  & \qw  & \targ\qwx[-3]\qw    & \qw                 & \ghost{3}\qw            & \multigate{1}{U_4} & \push{\dashedwire}\qw & \multigate{1}{U_4^{\dagger}} & \qw                          & \targ\qwx[-3]\qw              & \ghost{0}    & \lstick{\ket{0}     }    \\
  \lstick{\ket{0}     }  & \qw  & \qw                 & \targ\qwx[-2]\qw    & \qw                     & \ghost{4}          & \push{\dashedwire}\qw & \ghost{U_4^{\dagger}}        & \targ\qwx[-2]\qw             & \qw                           &  \qw            & \lstick{\ket{0}     }  \\
  \lstick{\ket{0}     }  & \qw  & \qw                 & \qw                 & \targ\qwx[-2]\qw        & \qw                & \push{\dashedwire}\qw & \qw                          & \qw                          & \qw                           & \ghost{y_1}  & \lstick{\ket{y_1}   }  \\      
  \lstick{\ket{0}     }  & \qw  & \qw                 & \qw                 & \qw                     & \targ\qwx[-2]\qw   & \ghost{5}             & \targ\qwx[-2]\qw             & \qw                          & \qw                           & \ghost{0}    & \lstick{\ket{0}     }    \\
  \lstick{\ket{0}     }  & \qw  & \qw                 & \qw                 & \qw                     & \qw                & \targ\qwx[-1]\qw      & \qw                          & \qw                          & \qw                           & \ghost{y_2}  & \lstick{\ket{y_2}     }  \\
}

%%     \end{center}
%%   \end{minipage}
  
%% \end{figure*}
%\vspace*{0.25cm}

\section{Proposed Method}
\label{Sec:proposed}
\subsection{Realizing Arbitrarily Controlled NOT Functions Up To A Relative Phase}
\label{Subsec:UNOT}

\begin{figure}[tbp]
    \centering
    \scalebox{0.8}{\Qcircuit @C=0.5em @R=0.2em @!R { \\
          \push{{x}_{0} }   & & \multigate{3}{x_k \cdot U}  & \qw   & \qw & &         & & \qw & \qw               & \qw                       & \qw      & \qw                       & \multigate{3}{U}              & \qw               & \qw      & \qw                         & \qw               & \qw & \\
          \push{{x}_{1} }   & & \ghost{x_k \cdot f}         & \qw   & \qw & &         & & \qw & \qw               & \qw                       & \qw      & \qw                       & \ghost{U}                      & \qw               & \qw      & \qw                         & \qw               & \qw & \\
          \push{\vdots}     & &                             &       &     & &\push{=} & &     &                   & \push{\vdots}             &          &  \push{\vdots}            &                                & \push{\vdots}     &          &                             & \push{\vdots}     &     & \\
          \push{{x}_{n-1} } & & \ghost{x_k \cdot f}         & \qw   & \qw & &         & & \qw & \qw               & \qw                       & \ctrl{1} & \qw                       & \ghost{U}                      & \qw               & \ctrl{1} & \qw                         & \qw               & \qw & \\
          \push{{x}_{n} }   & & \push{\rtoftarg}\qw\qwx[-1] & \qw   & \qw & &         & & \qw & \gate{\mathrm{H}} & \gate{\mathrm{T}}         & \targ    & \gate{\mathrm{T^\dagger}} & \targ\qwx[-1]    & \gate{\mathrm{T}} & \targ    & \gate{\mathrm{T^{\dagger}}} & \gate{\mathrm{H}} & \qw & \qw\\
}
                            
}
    \caption{A U-controlled RTOF gate}
    \label{fig-urtof}
\end{figure}

To get to a point where we can realize arbitrary Boolean functions up to a relative phase, we first observe that the RTOF depicted in Fig.~\ref{fig-rtof} can replace the center CNOT with a $n$-input Toffoli gate or U-controlled NOT (UCNOT) as depicted in Fig.~\ref{fig-urtof} which was observed by Maslov~\cite{bib-rtof-maslov}. However, as noted in~\cite{bib-amy-phase-state}, the introduction of a relative phase here has an effect on the final Boolean function as the final Hadamard gate rotates the relative phase back into state space. Therefore, we take the $C-X^*$ construction that is depicted in~\cite{bib-amy-phase-state}. We can generalize this $n$-input Toffoli gate to an arbitrarily controlled NOT gate as depicted in Fig.~\ref{fig-urtof}. 

\begin{figure}[tb]
    \centering
    \scalebox{0.8} {
    \Qcircuit @C=0.5em @R=0.2em @!R { \\
          \push{{x}_{0} }     & & \multigate{3}{D} & \multigate{3}{x_k \cdot g} & \qw   & \qw & &         & & \qw               & \qw                     & \multigate{3}{g} & \qw                     & \multigate{3}{f}  & \qw               & \multigate{3}{g} & \qw                        & \multigate{3}{f}   & \qw \\
          \push{{x}_{1} }     & & \ghost{D}        & \ghost{x_k \cdot g}        & \qw   & \qw & &         & & \qw               & \qw                     & \ghost{g}\qw     & \qw                     & \ghost{f}         & \qw               & \ghost{g}        & \qw                        & \ghost{f}          & \qw \\
          \push{\vdots}      & &                  &                            &       & \qw & &\push{=} & &                    & \push{\vdots}           &                  &  \push{\vdots}         &                   & \push{\vdots}     &                  &                            &                   & \push{\vdots} \\
          \push{{x}_{n} }     & & \ghost{D}        & \ghost{x_k \cdot g}       & \qw   & \qw & &         & & \qw               & \qw                     & \ghost{g}\qw     & \qw                     & \ghost{f}         & \qw               & \ghost{g}        & \qw                        & \ghost{f}          & \qw \\
          \push{{\ket{0}}   }& & \ghost{D}        & \targ\qwx[-1]              & \qw   & \qw & &         & & \gate{\mathrm{H}} & \gate{\mathrm{T^\dagger}} & \targ\qwx[-1]    & \gate{\mathrm{T^\dagger}} & \targ\qwx[-1]    & \gate{\mathrm{T}} & \targ\qwx[-1]    & \gate{\mathrm{T^{\dagger}}} & \qw                 &\gate{\mathrm{H}} \\
}


    }
    \caption{Generalization of Fig.~\ref{fig-urtof}, which also implements the Boolean product of two functions}
    \label{fig-f-dot-g}
\end{figure}

We can further generalize this. Recall from Section~\ref{Sec:Exist} that a Toffoli gate realizes a Boolean product of its inputs on its output. Just like above, we can use the same reasoning to generalize both inputs into arbitrary Boolean functions~\cite{bib-amy-phase-state} . This means that we can use the construction in Fig.~\ref{fig-f-dot-g} to realize a product between two arbitrary Boolean functions up to a relative phase. 
% \begin{figure*}[t]
%   \begin{minipage}{0.45\linewidth}
%     \begin{center}
%       \scalebox{0.5} {
%         \Qcircuit @C=0.5em @R=0.2em @!R { \\
          \push{{x}_{0} }     & & \multigate{5}{U_F}             & \qw  & \qw & &         & & \qw               & \qw               & \qw           & \qw                       & \multigate{5}{U_{F_{x_k=1}}}    & \qw               & \qw           & \qw                         & \multigate{5}{U_{F_{x_k=1}}^{\dagger}} & \qw           & \qw & \qw               & \qw             & \qw                       & \multigate{5}{U_{F_{x_k=0}}} & \qw               & \qw             & \qw                          & \multigate{5}{U_{F_{x_k=0}^{\dagger}}}    & \qw               & \qw\\
          \push{{x}_{1} }     & & \ghost{U_F}                    & \qw  & \qw & &         & & \qw               & \qw               & \qw           & \qw                       & \ghost{U_{F_{x_k=1}}}\qw        & \qw               & \qw           & \qw                         & \ghost{U_{F_{x_k=1}}^{\dagger}}        & \qw           & \qw & \qw               & \qw             & \qw                       & \ghost{U_{F_{x_k=0}}}        & \qw               & \qw             & \qw                          & \ghost{U_{F_{x_k=0}}^{\dagger}}\qw        & \qw               & \qw\\
          \push{\vdots}       & &                                &      & \qw & &         & &                   & \push{\vdots}     &               &  \push{\vdots}            &                                 & \push{\vdots}     &               &                             &                                        & \push{\vdots} &     &                   &                 &  \push{\vdots}            &                              & \push{\vdots}     &                 &                              &                                           & \push{\vdots}     &    \\
          \push{{x}_{k} }     & & \ghost{U_F}                    & \qw  & \qw & &         & & \qw               & \qw               & \ctrl{3}      & \qw                       & \ghost{U_{F_{x_k=1}}}\qw        & \qw               & \ctrl{3}      & \qw                         & \ghost{U_{F_{x_k=1}}^{\dagger}}        & \qw           & \qw & \gate{X}          & \ctrl{3}        & \qw                       & \ghost{U_{F_{x_k=0}}}\qw     & \qw               & \ctrl{3}        & \qw                          & \ghost{U_{F_{x_k=0}}^{\dagger}}\qw        & \gate{X}          & \qw\\
          \push{\vdots}       & &                                &      & \qw & &\push{=} & &                   & \push{\vdots}     &               & \push{\vdots}             &                                 & \push{\vdots}     &               &                             &                                        & \push{\vdots} &     &                   &                 &  \push{\vdots}            &                              & \push{\vdots}     &                 &                              &                                           & \push{\vdots}     &    \\
          \push{{x}_{n} }     & & \ghost{U_F}                    & \qw  & \qw & &         & & \qw               & \qw               & \qw           & \qw                       & \ghost{U_{F_{x_k=1}}}\qw        & \qw               & \qw           & \qw                         & \ghost{U_{F_{x_k=1}}^{\dagger}}        & \qw           & \qw & \qw               & \qw             & \qw                       & \ghost{U_{F_{x_k=0}}}\qw     & \qw               & \qw             & \qw                          & \ghost{U_{F_{x_k=0}}^{\dagger}}\qw        & \qw               & \qw\\
          \push{{\ket{0}}   } & & \push{\rtoftarg}\qw\qwx[-1]    & \qw  & \qw & &         & & \gate{\mathrm{H}} & \gate{\mathrm{T}} & \targ\qwx[-3] & \gate{\mathrm{T^\dagger}} & \push{\rtoftarg}\qw\qwx[-1]     & \gate{\mathrm{T}} & \targ\qwx[-3] & \gate{\mathrm{T^{\dagger}}} & \push{\rtoftarg}\qw\qwx[-1]            & \qw           & \qw & \gate{\mathrm{T}} & \targ\qwx[-1]   & \gate{\mathrm{T^\dagger}} & \push{\rtoftarg}\qw\qwx[-1]  & \gate{\mathrm{T}} & \targ\qwx[-1]   & \gate{\mathrm{T^{\dagger}}}  & \push{\rtoftarg}\qw\qwx[-1]               & \gate{\mathrm{H}} & \qw\\
}
%                              & &                                &      &     & &         & & \push{$g_0$}      & \push{$g_1$}      & \push{$g_2$}  & \push{$g_3$}              & \push{$g_4$}                & \push{$g_5$}      & \push{$g_6$}  & \push{$g_7$}                & \push{$g_8$}                 &               &     & \push{$g_9$}      & \push{$g_{10}$} & \push{$g_{11}$}           & \push{$g_{12}$}             & \push{$g_{13}$}   & \push{$g_{14}$} & \push{$g_{15}$}                     & \push{$g_{16}$}                    & \push{$g_{17}$}

%       }
%     \end{center}
%     \caption{Balanced Shannon's Decomposition}
%     \label{fig-shannon}
%   \end{minipage}
%   \hspace{0.5cm}
%   \begin{minipage}{0.45\linewidth}
%     \begin{center}
%       \scalebox{0.5} {
%         \Qcircuit @C=0.5em @R=0.2em @!R { \\
          \push{{x}_{0} }     & & \multigate{5}{F} & \qw  & \qw & &         & & \multigate{5}{F_{x_k=0}} & \qw               & \qw               & \qw            & \qw                       & \multigate{5}{F_{x_k=0} \oplus F_{x_k=1}} & \qw               & \qw           & \qw                          & \qw               & \qw\\
          \push{{x}_{1} }     & & \ghost{F}        & \qw  & \qw & &         & & \ghost{F_{x_k=0}}\qw     & \qw               & \qw               & \qw            & \qw                       & \ghost{F_{x_k=0} \oplus F_{x_k=1}}        & \qw               & \qw           & \qw                          & \qw               & \qw\\
          \push{\vdots}       & &                  &      & \qw & &         & &                          &                   &                   &                &  \push{\vdots}            &                                           & \push{\vdots}     &               &                              & \push{\vdots}     &    \\
          \push{{x}_{k} }     & & \ghost{F}        & \qw  & \qw & &         & & \ghost{F_{x_k=0}}\qw     & \qw               &                   & \ctrl{3}       & \qw                       & \ghost{F_{x_k=0} \oplus F_{x_k=1}}\qw     & \qw               & \ctrl{3}      & \qw                          &                   & \qw\\
          \push{\vdots}       & &                  &      & \qw & &\push{=} & &                          &                   &                   &                &  \push{\vdots}            &                                           & \push{\vdots}     &               &                              & \push{\vdots}     &    \\
          \push{{x}_{n} }     & & \ghost{F}        & \qw  & \qw & &         & & \ghost{F_{x_k=0}}\qw     & \qw               & \qw               & \qw            & \qw                       & \ghost{F_{x_k=0} \oplus F_{x_k=1}}\qw     & \qw               & \qw           & \qw                          & \qw               & \qw\\
          \push{{\ket{0}}   } & & \targ\qwx[-1]    & \qw  & \qw & &         & & \targ\qwx[-1]            & \gate{\mathrm{H}} & \gate{\mathrm{T}} & \targ\qwx[-1]  & \gate{\mathrm{T^\dagger}} & \targ\qwx[-1]                             & \gate{\mathrm{T}} & \targ\qwx[-1] & \gate{\mathrm{T^{\dagger}}}  & \gate{\mathrm{H}} & \qw\\
}

%       }
%     \end{center}
%     \caption{Unbalanced Positive Davio Decomposition}
%     \label{fig-davio-pos-unb}
%   \end{minipage}
% %\vspace*{0.5cm}
% \end{figure*}

\subsection{Using Shannon's Decomposition To Realize Arbitrarily Controlled Single Target Gates}
\label{Subsec:shannon}

\begin{figure}[tbp]
    \centering
      \scalebox{0.7} {
        \Qcircuit @C=0.5em @R=0.2em @!R { \\
          \push{{x}_{0} }  & \multigate{5}{x_k \cdot F} & \qw      & \multigate{5}{\bar{x_k} \cdot F'}&\qw\\
          \push{{x}_{1} }  & \ghost{x_k \cdot F}        & \qw      & \ghost{\bar{x_k} \cdot F'}       &\qw\\
          \push{\vdots}    &                          &          &                                &   \\
          \push{{x}_{k} }  & \ghost{x_k \cdot F}        & \gate{X} & \ghost{\bar{x_k} \cdot F'}       &\qw\\
          \push{\vdots}    &                          &          &                                &   \\
          \push{{x}_{n} }  & \ghost{x_k \cdot F}        & \qw      & \ghost{\bar{x_k} \cdot F'}       &\qw\\
          \push{{\ket{0}}} & \targ\qwx[-1]            & \qw      & \targ\qwx[-1]                  &\qw\\
}

      }
    \caption{A generalized way of conducting the OR for the Shannon decomposition}
    \label{fig-bool-or}
\end{figure}

Recall Shannon's decomposition $f(x_n,...,x_1) = x_n \cdot f(x_n=1,x_{n-1},\cdots,x_0) + \bar{x_n} \cdot f(x_n=0,x_{n-1},\cdots,x_0)$.
Using the constructions mentioned in Section~\ref{Subsec:UNOT}, we thus have a way to recursively construct
arbitrary Boolean functions up to a relative phase as a single target controlled NOT gate.

As a prelude to our discussion, we observe that $x \cdot a  + \bar{x} \cdot b = x \cdot a \oplus \bar{x} \cdot b$.
This means that we can implement this function as a series of Toffoli gates. Generalizing to UCNOT gates, we get
the construction shown in Fig.~\ref{fig-bool-or}, that we will be using to construct our implementations in the
next two sections.


First, from the Shannon's decomposition
formula~$f(x_n,...,x_1) = x_n \cdot f(x_n=1,x_{n-1},\cdots,x_0) + \bar{x_n} \cdot f(x_n=0,x_{n-1},\cdots,x_0)$
where we can take $g=f(x_n=1,x_{n-1},\cdots,x_0)$ and $g^{\prime}=f(x_n=0,x_{n-1},\cdots,x_0)$. As always with
Shannon's decomposition, we can decompose using an arbitrary variable order. (We will defer discussion of how we
determine this variable order to Section~\ref{Sec:exp}). We take Fig.~\ref{fig-f-dot-g} to insert our function, and
we get the construction as shown in Fig.~\ref{app-shannon}. We then take Fig.~\ref{fig-bool-or} and replace each
UCNOT in it with one of the terms in the Shannon's decomposition that we just made. As noted in
Section~\ref{Subsec:UNOT}, it is easy to see that this expression implements the Sum of Products (SOP) expression
of Shannon's decomposition. We then take each of $g$ and $g^{\prime}$ and repeat the process recursively until we
are left with a simple CNOT as our U-controlled NOT expression.

If it is ``False'', we have a False-controlled-NOT, meaning an identity function. This means that if one side of
the Shannon decomposition evaluates to either ``True'' or ``False'', we can merely replace it with a NOT gate or an
identity (meaning nothing at all), respectively, without expending any additional T-gates to implement the
function. As we will see in the next example, we can express the relative-phase Toffoli gate construction
in~\cite{bib-amy-phase-state} as a special case of this construction

Note that with this decomposition, when the function evaluates as ``True'', we end up with a True-controlled-NOT,
meaning, a NOT function. Because the Shannon decomposition now evaluates to either $x_k \cdot F_{x=1} + \bar{x}$ or
$x_k + \bar{x_k} \cdot F_{x=0}$, we can just implement the side with a True expression using an $x_k$ or
$\bar{x_k}$ controlled~CNOT gate. 

\begin{example}
  \label{ex-bool}
  We attempt to implement the function $F(x_3,x_2,x_1) = x_3x_2x_1 + \bar{x_3}(x_1 + x_2)$ as its relative phase equivalent.

  For the $F(x_3 = 1,x_2,x_1)$ case, we note that this construction just looks like a 3-input Toffoli gate. The first stage of this decomposition looks just like the
  decomposition in Fig~\ref{rtof}. We can see that if we
  continue Shannon's decomposition, one or the other side of the
  equation will always evaluate to 0. Therefore it is easy to see that it simplifies to the
  decomposition in Fig~\ref{fig-f-dot-g-in}, which is the decomposition in ~\ref{bib-amy-phase-state}.
  (As an aside, it generalizes it slightly as the algorithm also automatically handles the negated input
  case without any additional information). %The result is labeled with a curly brace in Fig.~\ref{fig-func-decomp}. 
  
  For the $F(x_3 = 1,x_2,x_1)$ case, we now have something that is not a straightforward sum of products term.
  We now move on to the next level (Note: The variable order here, $x_3 \rightarrow x_2 \rightarrow x_1$, is
  chosen to best illustrate the working of the method. We treat optimal variable order later). We see that
  $F(x_3 = 0,x_2=0,x_1) = x_1$. Following the algorithm here results in something that looks like an $x_2$
  and $x_1$ controlled Toffoli gate. We then move on to $F(x_3 = 0,x_2=1,x_1)$ and find that it is $True$.
  This is one of the special cases in our algorithm, and so we can implement $x_2 \cdot F(x_3 = 0,x_2=1,x_1)$
  as an $x_2$ controlled NOT gate. We then use the construction depicted in Fig~\ref{fig-bool-or} to compose
  the two functions. %The result is depicted Fig.~\ref{fig-func-decomp}, indicated by the dashed box. 
  
  Finally we compose $F(x_3 = 1,x_2,x_1)$ and $F(x_3 = 1,x_2,x_1)$ as in Fig~\ref{fig-bool-or} again. The result is depicted in Fig.~\ref{app-func-decomp} in ~\ref{App:FuncDecomp}.
\end{example}

However, as we will see later in Section \ref{Sec:Exp}, some expressions are not amenable to a purely Shannon
Decomposition. Several of them have highly complex (many BDD nodes) representation in the Shannon Decomposition.
We must then introduce a different representation to take them into account.

\subsection{Integrating Davio Expansions}
\label{Subsec:davio}
As can be seen in the previous section, a significant portion of the cost comes from multiplication.
Therefore it is useful to consider a method that takes advantage of {\it Exclusive Sum of Products} (ESOP)
expressions. Recall the positive Davio Expansion $F=F_1\oplus x_n \cdot F_2$ and the corresponding
negative Davio expansion $F=F_0 \oplus \bar{x_n} \cdot F_2$ where $F_0 = F(x_n=0)$,$F_1 = F(x_n=1)$,
and $F_2 = F_0 \oplus F_1$. These expansions have the advantage that at every level, only one
multiplication takes place, which already reduces the maximum possible T-gate use at any node from
8 to 4. 

\subsection{Using Balanced and Unbalanced Relative Phase Functions}
\label{Subsec:balanced}
As Amy et al. ~\cite{bib-amy-phase-state} pointed out, because any multiplication is first ``constructed''
in phase space and then rotated back into state space with a Hadamard gate, adding a relative phase circuit
inside that logic with a \emph{target dependent} phase (meaning one whose value depends on the target bit
of the function) in an \emph{unbalanced} construction, such as in Fig~\ref{fig-urtof}, will see its relative
phase added on to the main Boolean function as a residue i.e. it will take the form of $F=x \cdot F_k + F_{k-1}$.
To keep this from happening, a \emph{balanced} construction such as in Fig~\ref{fig-f-dot-g} is used, with pairs
of relative phase circuits and their inverses to cancel out the target dependence that the Hadamard gate on the
bit will act on. In the process, such a construction creates a function with a \emph{target independent}relative
phase. This means that it can be used inside an unbalanced construction without incurring residue since the
Hadamard gate now does not have any dependence to act on.

Amy et al. propose a method where balanced and unbalanced constructions are alternated at different levels of the
construction of an MCT to reduce the amount of times a function has to be doubled to cancel itself out. We can
easily extend this to our method of constructing arbitrary Boolean functions above. We start at the root node
with an unbalanced construction, and any children that are constructed recursively will have to be constructed
in a balanced manner. The children of those in turn can be unbalanced, and so on. 

In doing so, we can reduce the dependency from $2^n$ to $2^{\frac{n}{2}}$. While still exponential, this has quite
a significant reduction in the regime we are working with. For example for $n=6$, the circuit is at MOST 1/8th of
the original cost! Merely introducing this methodology results in significant savings for larger LUT-sizes, though
its effect could be diminished for smaller ones.

This has additional ramifications when combined with the Davio Expansions. Since there is only one multiplication
performed, the non-multiplied term ($F_0$ in the positive Davio case and $F_1$ in the negative Davio case) can
be implemented as an \emph{unbalanced} relative phase circuit, provided the parent construction can also be
unbalanced. This means there is no need to double the logic inside to cancel the residue, and so the exponential
is further reduced. This can reduce the exponential relation even further.

\subsection{Relative Phase Synthesis Algorithm}
\label{Subsec:Pro:Improved}
We now try to integrate the Sec~\ref{Subsec:shannon}, Sec~\ref{Subsec:balanced}, and ~\ref{Subsec:davio}.

As a prelude to our discussion, we observe that $x \cdot a  + \bar{x} \cdot b = x \cdot a \oplus \bar{x} \cdot b$.
This means that we can implement this function as a series of Toffoli gates. Generalizing to UCNOT gates, we get
the construction shown in Fig.~\ref{fig-bool-or}.

First, we determine the variable order to decompose the function in. As the details of the variable order
computation are not important to the algorithm itself, we will defer discussion of how we determine
this variable order to our discussion of the implementation in Section~\ref{Sec:exp}.

If the function is neither ``True'' nor ``False'', the function is expanded using the Shannon, Positive
Davio, or Negative Davio expansions, independently. We go over constructions for them in Appendix
~\ref{App:ManyDecomp}. The constituent subfunctions of each expansion are then processed in the same way by
recursing the function with the decomposed Boolean function as input.. When the logic can be implemented as an
unbalanced function, such as in the unmultiplied term of the Davio expansion or as balanced logic gates, that is
relayed to the next level of the recursion via an input.

If it is ``False'', we have a False-controlled-NOT, meaning an identity function. In this case we remove the term
with this function entirely.

Note that with this decomposition, when the function evaluates as ``True'', we end up with a True-controlled-NOT, meaning, a NOT function. With the multiplied terms such as the $x_k \cdot F_{x=1}$ in Shannon's Decomposition or $\bar{x_k} \cdot (F_{x=0} \oplus F_{x=1})$ in the Davio Expansion, we simply get $x_k$ or $\bar{x_k}$ respectively. This means we can just implement the side with a True expression using an $x_k$ or $\bar{x_k}$ controlled~CNOT gate. 

As the recursion terminates, the cost of each of the Shannon, Positive Davio, or Negative Davio expansions is
compared with each other, and the minimum is chosen. The result is then passed on to the recursion level above it.
The process is then repeated with all of the variable orders and the minimum is picked from the results. The
resulting algorithm is described in Alg.~\ref{algo-new}


\begin{algorithm}[tb]
  \caption{relative\_phase(BF: a Boolean function, balanced: Boolean indicating whether this function is balanced)}
  \label{algo-new}
  \begin{algorithmic}
%    \State BF $\leftarrow$ A Boolean Function
    \State QC $\leftarrow$ an empty quantum circuit with registers
    \State Output $\leftarrow$ the output bit of the Boolean Function
    \State order $\leftarrow$ BF.find\_optimal\_order()
    \If{$BF == True$}
    \State return CNOT
    \ElsIf{$BF == False$}
    \State return QC $\leftarrow$ Identity(Output)
    \ElsIf{$BF$ is a leaf node}
    \State return CNOT
    \Else
    \State    //The Shannon $F_0$ and $F_1$ alternate balanced and unbalanced
    \State    $F_0$ $\leftarrow$ relative\_phase($BF_{x_{order[first]}\leftarrow 0)}$,!balanced) 
    \State    $F_1$ $\leftarrow$ relative\_phase($BF_{x_{order[first]}\leftarrow 1}$,!balanced) 
    \State
    \State    //The Davio $F_0$ and $F_1$ inherit the parent's balanced attribute
    \State    $(F_0)_{dav}$ $\leftarrow$ relative\_phase($BF_{x_{order[first]}\leftarrow 0)}$,balanced) 
    \State    $(F_1)_{dav}$ $\leftarrow$ relative\_phase($BF_{x_{order[first]}\leftarrow 1}$,balanced) 
    \State
    \State    $F_2$ $\leftarrow$ relative\_phase($BF_{x_{order[first]}\leftarrow 1} \oplus BF_{x_{order[first]}\leftarrow 0}$,
    \State                       !balanced)
    \State
    \State    $shannon$ $\leftarrow$ compose $F_0$ and $F_1$ in a balanced manner if balanced or unbalanced if not.
    \State    $davio_{pos}$ $\leftarrow$ compose $(F_0)_{dav}$ and $F_2$ in a balanced manner if balanced is true or unbalanced if not.
    \State    $davio_{neg}$ $\leftarrow$ compose $(F_1)_{dav}$ and $F_2$ in a balanced manner if balanced is true or unbalanced if not.
    \State    return one of $[shannon, davio_{pos},davio_{neg}]$ with the min T-count
    \EndIf
    \State Repeat for every variable order and pick the minimum
  \end{algorithmic}
  
\end{algorithm}


\begin{example}
  \label{ex-davio}
    We again attempt to implement the function $F(x_3,x_2,x_1) = x_3x_2x_1 + \bar{x_3}(x_1 + x_2)$ as its relative phase equivalent.
    Just as before, we first take the positive Davio expansion about $x_3$ (for simplicity, we will only follow the Davio branches
  in this example. However, as  we will see later, they turn out to be optimal for this function). We get $F_0 = x_1 + x_2$ and
  $F_2 = F_0 \oplus F_1 = (x_1 + x_2) \oplus x_1 \cdot x_2 = x_1 \oplus x_2$. Note that for this level we can tolerate an unbalanced
  function, so for the $F_0$ branch we set the ``balanced'' parameter to false and on the $F_2$ branch we set it to true.

  We then proceed to $x_2$. We decompose $F_0 = (F_0)_0 \oplus x_2 \cdot ((F_0)_2)$. Since $(F_0)_0 = x_1$, we just implement it as an $x_1$
  controlled CNOT. After calculating $(F_0)_1 = x_1 + 1 = 1$, we get $(F_0)_2 = (F_0)_0 \oplus (F_0)_1 = x_1 \oplus 1 = \bar{x_1}$
  Thus $x_2 \cdot (F_0)_2$ is implemented according to a balanced function multiplication such as in Fig.~\ref{fig-f-dot-g}. We then compose
  these similarly to Fig.~\ref{app-davio-pos-unb} but with a balanced constructionand terminate the recursion. We can see at this point that
  there are 4 T-gates.

  We now go to the $F_2$ branch. Here, since $x_1 \oplus x_2$ can be implemented using an $x_1$ controlled CNOT and $x_2$ controlled CNOT
  in series, we synthesize that and terminate the recursion.

  Back at the top level, we compose the subcircuits for $F_0$ and $F_2$ as a negative unbalanced Davio construction, adding an additional 4
  T-gates. This brings us to a total of 8 T-gates. Compare this to the 24 T-gates of our original methodology, 1/3 the size. We omit
  the computation of the other branches and variable orders in the interest of space, but the reader can perform them  quickly to
  see that the branch we followed is the optimal branch (or one of them. In this function there are multiple because $x_1$  and
  $x_2$ are symmetric, and several Davio branches have similar costs). 
  The resultant circuit is depicted in Fig.~\ref{fig-func-decomp-davio} in \ref{App:FuncDecomp}

\end{example}

\begin{example}

%% \begin{algorithm}[tbp]
%%   \caption{shannon\_decomp(BF: a Boolean function)}
%%   \label{algo-shannon}
%%   \begin{algorithmic}
%% %    \State BF $\leftarrow$ A Boolean Function
%%     \State QC $\leftarrow$ an empty quantum circuit with registers
%%     \State Output $\leftarrow$ the output bit of the Boolean Function
%%     \State order $\leftarrow$ BF.find\_optimal\_order()
%%     \If{$BF == True$}
%%     \State return CNOT
%%     \ElsIf{$BF == False$}
%%     \State return QC $\leftarrow$ Identity(Output)
%%     \ElsIf{$BF$ is a leaf node}
%%     \State return CNOT
%%     \Else
%%     \State $F_0$ $\leftarrow$ shannon\_decomp($BF_{x_{order[first]}\leftarrow 0)}$)
%%     \State $F_1$ $\leftarrow$ shannon\_decomp($BF_{x_{order[first]}\leftarrow 1}$)
%%     \State return a UCNOT implemented by composing $F_0$ + $F_1$ in the manner shown in Fig.~\ref{app-shannon}
%%     \EndIf
%%   \end{algorithmic}
  
%% \end{algorithm}

\begin{algorithm}[tb]
  \caption{lut\_synthesis}
  \label{algo-lut}
  \begin{algorithmic}
    \State lut\_net $\leftarrow$ a LUT network
    \State esop (Boolean Function) $\leftarrow$ a method to generate
    an exact QC to calculate a Boolean function implemented as a UCNOT gate 
    \For{every node in lut\_net}
    \If{node is intermediate node}
    \State QC $\leftarrow$ add relative\_phase(node's Boolean function) on a dirty ancilla
    \State QC $\leftarrow$ add inverse of relative\_phase(node's 
    Boolean function) where we restore the state of the same dirty ancilla
    \ElsIf{node is output node}
    \State QC $\leftarrow$ add esop (node's Boolean function) acting on a dirty ancilla
    \EndIf
    \EndFor
  \end{algorithmic}
%\vspace*{-0.25cm}
\end{algorithm}

\subsection{Synthesizing A Boolean Function Using LUT-networks and Shannon Decomposition RTOF}

Using Algorithm~\ref{algo-new} we are now ready to propose our LUT-based synthesis method 

First we compose the circuit as in Sec.~\ref{Sec:lhrs}. However, this time, we replace the alternating pairs of compute/uncompute gates with the respective relative phase version using Algorithm~\ref{algo-shannon} as well as its associated inverse. If there is no matching uncompute logic i.e., a node acting on an output qubit, we generate the logic exactly. The method is enumerated explicitly in Algorithm~\ref{algo-lut}.

This QC is then parsed by our algorithm, synthesizing each UCNOT into relative phase versions of the partial Boolean functions in each intermediate node (as well as their inverse), and exact versions in each output node. These are then recombined into a quantum circuit. 

\label{ex-lut-rp}
  We follow the steps in Ex.~\ref{ex-lut}. We then parse through LUTs
  $1$,$2$,$1^{\dagger}$,$2^{\dagger}$,$4$,and $4^{\dagger}$ and
  generate their relative phase versions. For nodes 3 and 5, we use an
  ESOP-based exact synthesis. 
  The result is depicted in Fig.~\ref{fig-lut-qc}.
\end{example}



\begin{table*}[ht]
  \begin{center}
    \caption{Experimental Results}
    \label{table-results}
    \scalebox{0.90}{
      \begin{tabular}{c|c|c|c|c|c|c|c|c|c|c|c|c}
        \hline
          Circuit             & \multicolumn{4}{c}{ESOP-based method}  &  \multicolumn{4}{|c|}{Proposed method} &  \multicolumn{4}{|c}{Previous method} \\ \cline{2-13}
                              & LUT6 (T-count) & LUT5 & LUT4 & LUT3    & LUT6 & LUT5 & LUT4 & LUT3              & LUT6 & LUT5  & LUT4 & LUT3 \\ \hline                              
        priority\_size\_2022  & 98664          & 0.56 & 0.32 & 0.26    & 0.14 & 0.08 & 0.06 & 0.05              & 0.27 & 0.51 & 0.19 & 0.16  \\ \hline            
        int2float\_size\_2022 & 22592          & 1.14 & 0.49 & 0.38    & 0.27 & 0.17 & 0.1  & 0.06              & 0.94 & 0.94 & 0.4 & 0.35  \\ \hline             
        i2c\_size\_2022       & 89288          & 0.69 & 0.4 & 0.27     & 0.61 & 0.27 & 0.18 & 0.13              & 0.94 & 0.68 & 0.31 & 0.23 \\ \hline             
        dec\_size\_2018       & 57344          & 0.5 & 0.41 & 0.2      & 1 & 0.39 & 0.3 & 0.15                  & 1 & 0.41 & 0.31 & 0.16 \\ \hline                
        ctrl\_size\_2022      & 7560           & 1.17 & 0.84 & 0.68    & 0.99 & 0.88 & 0.35 & 0.24              & 0.99 & 1.12 & 0.51 & 0.41\\ \hline              
        cavlc\_size\_2022     & 116896         & 0.8 & 0.49 & 0.41     & 0.13 & 0.1 & 0.07 & 0.06               & 0.63 & 0.6 & 0.3 & 0.29 \\ \hline               
        bar\_size\_2015       & 176400         & 0.94 & 0.03 & 0.03    & 0.15 & 0.21 & 0.07 & 0.07              & 7.15 & 1.4 & 0.21 & 0.21\\ \hline               
        arbiter\_size\_2022   & 175640         & 0.68 & 0.42 & 0.28    & 0.21 & 0.13 & 0.11 & 0.05              & 0.35 & 0.63 & 0.27 & 0.17 \\ \hline             
        adder\_size\_2022     & 4992           & 9.75 & 6.38 & 1.94    & 1 & 1.91 & 1.95 & 1.66                 & 1 & 7.16 & 4.06 & 2.54  \\ \hline               
                              &  \multicolumn{4}{c|}{CPU time of the proposed method (sec)}  \\ \cline{1-5}
                              & LUT6 & LUT5  & LUT4 & LUT3 \\ \cline{1-5}
        priority\_size\_2022  & 14.97 & 8.07 & 5.01 & 0.79 \\ \cline{1-5}
        int2float\_size\_2022 & 13.31 & 4.02 & 6.03 & 0.87 \\ \cline{1-5}
        i2c\_size\_2022       & 115.00 & 33.35 & 19.73 & 12.36 \\ \cline{1-5}
        dec\_size\_2018       & 12.23 & 7.03 & 18.72 & 21.09 \\ \cline{1-5}
        ctrl\_size\_2022      & 10.17 & 5.57 & 3.21 & 1.11 \\ \cline{1-5}
        cavlc\_size\_2022     & 77.91 & 31.92 & 2.27 & 2.23 \\ \cline{1-5}
        bar\_size\_2015       & 410.26 & 215.09 & 29.92 & 35.28 \\ \cline{1-5}
        arbiter\_size\_2022   & 53.63 & 49.48 & 10.00 & 21.09 \\ \cline{1-5}
        adder\_size\_2022     & 144.37 & 98.02 & 40.37 & 33.56 \\\cline{1-5}
      \end{tabular}
   }
  \end{center}
%\vspace*{0.25cm}
\end{table*}











%\vspace*{0.5cm}
\section{Experimental Results}
\label{Sec:exp}
\subsection{Methodology}

Boolean functions are taken from the EPFL Combinational Benchmark Suite~\cite{bib-epfl} and generated using the ABC tool~\cite{bib-abc}. We choose LUT-sizes of 3, 4, 5, and 6, variables. as in~\cite{bib-barenco-elementary}, this is theorized as the upper bound to the tractable implementations of Boolean functions in the quantum space. This LUT network is converted to a quantum circuit implemented in Qiskit~\cite{bib-qiskit}, including any uncomputation logic. 

%The Shannon Decomposition algorithm uses the Tulip Control BDD library~\cite{bib-tulip-dd} to implement its Boolean functions, as well as reorder the variables for optimality (reordering is done using Rudell's Sifting Algorithm~\cite{bib-rudell}).Any exact synthesis of Boolean functions is conducted using Qiskit's ClassicalFunction functionality, which uses the ESOP-based PKRM algorithm~\cite{bib-pkrm}.
The Decomposition algorithm is implemented in Python. In addition to the algorithm described in Algorithm~\ref{algo-imp}, we also compare several modified versions. A version that only utilizes the balanced and unbalanced constructions(noDavio) and one that only utilizes the Davio expansion (noBalance) are compared to the original to measure the contribution of each of these respective elements. Any exact synthesis of Boolean functions is conducted using Qiskit's ClassicalFunction functionality, which uses the ESOP-based PKRM algorithm~\cite{bib-pkrm}.

These results are then compared to a naive implementation of the LUT-based synthesis method, wherein each LUT is synthesized as an exact (non relative phase) implementation of their Boolean functions, as well as the original Shannon's Decomposition based algorithm. Because the relative phase functions are constructed without any ancilla, we assume that there are plenty of bits to be reused at any time, so any MPMCT (Multi-Polar Multiply-Controlled Toffoli) gates are synthesized directly using Gray-code sequence of operations for $n<5$, or decomposed according to the Barenco dirty ancilla method~\cite{bib-barenco-elementary}. This same decomposition method is also used in the proposed method for any output node that requires an exact implementation.

As in~\cite{bib-barenco-elementary}, we note that for values of $n<7$, a direct implementation of the Toffoli gate (i.e. using a Gray-code sequence of operations) the number of gates to implement each MCT gate scales with $2^{n-1}$.

Our experimental results are displayed in
Table~\ref{table-results}. We compare the ratio of the T-count for
each mapped LUT-size generated by either the ESOP-based method or the proposed 
method to that of the LUT6 in the second column generated by the
ESOP-based method as a baseline.  
%In addition to comparing the numbers for each LUT map size
%with the proposed method to this baseline, we also include the numbers
%for their counterparts generated by the ESOP method, to measure the
%overall effect of decreasing the LUT-size on the T-count.
Recall that in the previous paragraph, we noted that the cost of each relative
phase LUT scales with $2^n$. This measures the contribution of
decreasing the LUT-size at the cost of increasing the number of gates 
and ancilla used. Finally, we also measure the execution time in
seconds of our proposed method of for each LUT-size. 

%Our experimental results in Table~\ref{table-results} use three sets of Numbers: vsLUTXESOP is the ratio of the difference between the T-count from using the proposed method and the ESOP method for each maximum LUT-size (LUTX), and that from using the ESOP method for that LUT-size i.e. $Tcount(LUTX_{ESOP}) - Tcount(LUTX_{prop}) / Tcount(LUTX_{ESOP})$ (displayed as \%). It compares the T-count of the proposed method for each X-input LUT size with synthesizing each LUT using ESOP based methods. This measures the effectiveness of the Shannon Decomposition Method in replacing exact logic. We also display a similar measure vsLUT6ESOP, which takes of the T-count of the proposed method for each size and measures the ratio of its difference with the T-count of the ESOP method for the 6-input maximum LUT case (LUT6) vs the T-count of that ESOP method i.e. $Tcount(LUT6_{ESOP}) - Tcount(LUTX_{prop}) / Tcount(LUT6_{ESOP})$ (displayed as \%). Because the Shannon decomposition method scales with $2^n$, comparing to the LUT6 ESOP numbers gives us a measure of the overall effectiveness of the algorithm at different LUT-sizes, integrating effects that the LUT-size might have on the overall cost of the circuit. Finally, we also measure the execution time ExecTime in seconds of each method. 

\subsection{Analysis}
In Table~\ref{table-results}, the best result for each circuit is written in bold; we see that in most cases the proposed method with LUT-size is three is the best. However, there remain cases where the proposed method underperforms the ESOP implementation, such as bar\_size\_2015. 
Qualitatively, a common trend among circuits that seem to have similar 
numbers between the proposed and ESOP-based methods is that they have large
numbers of highly complex, maximum width gates. This means more
complex BDDs, and because the cost of the Shannon decomposition scales 
exponentially with the number of nodes, this will close the gap with 
the ESOP-based method. This is the case with bar\_size\_2015, where
the vast majority of its gates are at the maximum LUT-size, leading to
losing out to ESOP across the board. Others, such as dec\_size\_2018,
instead have maximum width but minimum complexity (i.e., a single MCT gate)
Boolean functions. As our proposed method may perform worse than
synthesizing these gates directly using Gray-code, this will lead to a
diminished advantage. Finally, a larger share of output nodes to 
intermediate nodes also contributes, as the method is identical 
between the two. A particularly interesting case is 
adder\_size\_2022. At LUT6, this circuit's Boolean logic network can
fit all of the logic for each bit into a single node, so there is no 
need to generate any intermediate nodes; therefore, every single LUT 
is generated identically to the ESOP-based method i.e., a ratio of
1. However, moving onto LUT5, we see an immediate jump in the numbers, 
ostensibly because of the $2^n$ proportionality. We thereafter see a 
drop in the ratio, though it still underperforms the ESOP-based method
even when LUT-size is 3. 


The balance between all of these elements seems to be the key in maximizing the advantage the proposed method has over direct synthesis methods. An interesting future topic of research would be to find how
to produce LUT networks that maximize this advantage. Looking further, it is apparent that at smaller LUT-sizes, the Shannon decomposition method produces a much clearer advantage over the direct method. Additionally, we see using both the Shannon decomposition method in conjunction with smaller LUT-sizes yields a clear advantage in almost all of the cases we see. While this yields an advantage in T-count, keep in mind that the required ancilla increases. Even if we use dirty ancilla, trading T-gates for logical qubits, enough dirty ancilla will result in congestion in a real system, leading to increased gate counts that could effectively trade T-gates for CNOTs, and thus T-gates for an increased error-rate. A future topic of research could possibly look into an implementation hardware aware version of this algorithm that considers these resources at the time of synthesis of the LUT-network.





%% \begin{table*}[t]
%%   \large
%%    \begin{center}
%%      \caption{LUT6 Experimental Results}
%%      \scalebox{1.0}{
%%        \begin{tabular}{|c|c|c|c|c|c|}
%%          \hline
%%          Circuit & Proposed & ESOP & TotGates & NumOuts & AvgNodes \\
%%          \hline
%%          priority\_size\_2022 & 27160 & 89800 & 180 & 8 & 3 \\
%%          \hline
%%          int2float\_size\_2022 & 21776 & 25528 & 33 & 7 & 6 \\
%%          \hline
%%          i2c\_size\_2022 & 89440 & 93584 & 251 & 141 & 5 \\
%%          \hline
%%          dec\_size\_2018 & 63288 & 112224 & 517 & 11 & 4 \\
%%          \hline
%%          ctrl\_size\_2022 & 12096 & 17712 & 48 & 4 & 3 \\
%%          \hline
%%          cavlc\_size\_2022 & 74064 & 113616 & 97 & 11 & 6 \\
%%          \hline
%%          bar\_size\_2015 & 1265472 & 180000 & 1018 & 6 & 4 \\
%%          \hline
%%          arbiter\_size\_2022 & 61328 & 176728 & 417 & 129 & 4 \\
%%          \hline
%%          adder\_size\_2022 & 4912 & 4912 & 129 & 129 & 7 \\
%%          \hline
%%        \end{tabular}
%%      }
%%      \label{table-results6}
%%      \end{center}
%% \vspace*{0.25cm}
%%  \end{table*}

%% \begin{table*}[t]
%%   \large
%%    \begin{center}
%%      \caption{LUT3 Experimental Results}
%%      \scalebox{1.0}{
%%        \begin{tabular}{|c|c|c|c|c|c|}
%%          \hline
%%          Circuit & Proposed & ESOP & TotGates & NumOuts & AvgNodes \\
%%          \hline
%%          priority\_size\_2022 & 15960 & 25328 & 400 & 8 & 2 \\
%%          \hline
%%          int2float\_size\_2022 & 7984 & 8480 & 149 & 7 & 2 \\
%%          \hline
%%          i2c\_size\_2022 & 22656 & 33088 & 709 & 11 & 2 \\
%%          \hline
%%          dec\_size\_2018 & 5312 & 19376 & 647 & 11 & 2 \\
%%          \hline
%%          ctrl\_size\_2022 & 3216 & 6544 & 104 & 4 & 2 \\
%%          \hline
%%          cavlc\_size\_2022 & 33904 & 47936 & 560 & 4 & 2 \\
%%          \hline
%%          bar\_size\_2015 & 42752 & 6176 & 1786 & 6 & 2 \\
%%          \hline
%%          arbiter\_size\_2022 & 30096 & 50080 & 1067 & 7 & 2 \\
%%          \hline
%%          adder\_size\_2022 & 16424 & 14912 & 506 & 12 & 2 \\
%%          \hline
%%        \end{tabular}

%%      }
%%      \label{table-results3}
%%      \end{center}
%% \vspace*{0.25cm}
%%  \end{table*}

%% \begin{table*}[t]
%%   \large
%%    \begin{center}
%%      \caption{LUT4 Experimental Results}
%%      \scalebox{1.0}{
%%        \begin{tabular}{|c|c|c|c|c|c|}
%%          \hline
%%          Circuit & Proposed & ESOP & TotGates & NumOuts & AvgNodes \\
%%          \hline
%%          priority\_size\_2022 & 19064 & 32928 & 378 & 8 & 2 \\
%%          \hline
%%          int2float\_size\_2022 & 8976 & 10760 & 139 & 7 & 2 \\
%%          \hline
%%          i2c\_size\_2022 & 31424 & 47856 & 615 & 11 & 2 \\
%%          \hline
%%          dec\_size\_2018 & 10344 & 38416 & 605 & 11 & 2 \\
%%          \hline
%%          ctrl\_size\_2022 & 4064 & 7640 & 96 & 4 & 2 \\
%%          \hline
%%          cavlc\_size\_2022 & 35104 & 56488 & 542 & 4 & 2 \\
%%          \hline
%%          bar\_size\_2015 & 42736 & 6032 & 1786 & 6 & 2 \\
%%          \hline
%%          arbiter\_size\_2022 & 50776 & 84944 & 935 & 7 & 2 \\
%%          \hline
%%          adder\_size\_2022 & 39352 & 37104 & 624 & 12 & 2 \\
%%          \hline
%%        \end{tabular}
%%      }
%%      \label{table-results4}
%%      \end{center}
%% \vspace*{0.25cm}
%%  \end{table*}

%% \begin{table*}[t]
%%   \large
%%    \begin{center}
%%      \caption{LUT5 Experimental Results}
%%      \scalebox{1.0}{
%%        \begin{tabular}{|c|c|c|c|c|c|}
%%          \hline
%%          Circuit & Proposed & ESOP & TotGates & NumOuts & AvgNodes \\
%%          \hline
%%          priority\_size\_2022 & 50800 & 52560 & 330 & 8 & 2 \\
%%          \hline
%%          int2float\_size\_2022 & 21576 & 27272 & 121 & 7 & 3 \\
%%          \hline
%%          i2c\_size\_2022 & 73960 & 87088 & 503 & 11 & 3 \\
%%          \hline
%%          dec\_size\_2018 & 15312 & 48272 & 591 & 11 & 2 \\
%%          \hline
%%          ctrl\_size\_2022 & 10136 & 13832 & 52 & 4 & 3 \\
%%          \hline
%%          cavlc\_size\_2022 & 71272 & 102504 & 372 & 4 & 3 \\
%%          \hline
%%          bar\_size\_2015 & 329360 & 178736 & 1692 & 6 & 3 \\
%%          \hline
%%          arbiter\_size\_2022 & 118480 & 128336 & 861 & 7 & 2 \\
%%          \hline
%%          adder\_size\_2022 & 114544 & 54128 & 378 & 12 & 3 \\
%%          \hline
%%        \end{tabular}

%%      }
%%      \label{table-results4}
%%      \end{center}
%% \vspace*{0.25cm}
%%  \end{table*}


\vspace*{-0.25cm}
\section{Conclusion and Future Research}
\label{Sec:concl}

This paper proposed the use of Davio expansions and unbalanced relative phase constructions to optimize relative phase constructions in their usage in LUT-based quantum circuit synthesis. These were meant to improve the $2^n$ proportionality of an earlier, Shannon Decomposition based method. The experimental results showed a clear advantage over this previous method, as well as a significant advantage over the naive method. There remains one interesting exception where the proposed method underperforms the naive method, however, and it relates to a special form of LUT network.

The results of this paper show a clear advantage to using our proposed method, as well as possible further avenues of inquiry for related research. As the experimental results show, there remain many variables to control to optimize in this method, and researching those further could prove fruitful.
%This paper proposed the use of relative phase Boolean functions to optimize LUT-based quantum circuit synthesis while maintaining a minimum of ancillae. While the Shannon decomposition allows us to the cost remains exponential with respect to the number of LUT inputs, limiting the LUT-size to under 6 inputs allows us to see a significant T-count advantage in the proposed method using relative phase Boolean functions.

%There still remain some cases where it underperforms the naive case so further research will be needed there. Furthermore, it remains an open question whether we can get better than exponential cost for the LUTs. As shown in~\cite{bib-lhrs-soeken}, increasing LUT size has possible advantages in gate count and depth. Furthermore, our research has shown that decreasing LUT-size has an advantage in T-count. Existing research in this topic has now uncovered several variables with which to toggle to minimize the T-count, and points to potential further research in the topic.


\bibliographystyle{ieicetr}%
\bibliography{ref}%

\vspace{-22mm}

\profile[img/clarino.pdf]{David Clarino}{received his BS EECS and BA Physics from UC Berkeley and MS from Santa Clara University. He is currently a student at the Graduate School of Information Science and Engineering, Ritsumeikan University, Shiga, Japan. His research interests include quantum circuit design and quantum circuit verification and he has previously worked on SoC design verification}

\vspace{-5mm}

\profile[accel.jpg]{Naoya Asada}{received his bachelor's degree from
College of Information Science and Engineering, Ritsumeikan
University in 2021. He received his master's degree from
Graduate School of Information Science and Engineering, Ritsumeikan
University in 2023. He is currently with Mitsubishi Electric Corp. 
}

\vspace{-5mm}

\profile[img/Yamashita.pdf]{Shigeru Yamashita}{is a professor at the
Department of Computer Science, College
of Information Science and Engineering,
Ritsumeikan University. He received
his B.E., M.E. and Ph.D. degrees in Information
Science from Kyoto University,
Kyoto, Japan, in 1993, 1995 and 2001, respectively.
His research interests include
new types of computation and logic synthesis for them. He received
the 2000 IEEE Circuits and Systems Society Transactions
on Computer-Aided Design of Integrated Circuits and Systems
Best Paper Award, SASIMI 2010 Best Paper Award, 2010 IPSJ
Yamashita SIG Research Award, and 2010 Marubun Academic
Achievement Award of the Marubun Research Promotion Foundation.
He is a senior member of IEEE, and a member of IPSJ.}
\end{document}
